\documentclass[../../../include/open-logic-section]{subfiles}

\begin{document}

\olfileid{sth}{z}{milestone}
\olsection{$\Z^-$: ఒక మైలురాయి}

తరువాతి విభాగంలో
\stagesinf{}ను మళ్లీ
పరిశీలిస్తాం. అయితే
అనంతత్వ స్వీకృతంతో
ఒక ముఖ్యమైన మైలురాయిని
చేరుకున్నాం. ఇప్పుడు
$\Zminus$ సిద్ధాంతానికి
అవసరమైన స్వీకృతాలన్నీ
మన దగ్గర ఉన్నాయి.
వివరంగా:

\begin{defn}
$\Zminus$ సిద్ధాంతంలో ఈ
స్వీకృతాలు ఉంటాయి:
సమితుల సమానత్వ సూత్రం,
సమ్మేళనం, జంటలు,
ఘాత సమితులు, అనంతత్వం,
వేరుచేయడం పథకంలోని
అన్ని సందర్భాలు.
\end{defn}

ఈ పేరు \emph{సెర్మెలో}
సమితి సిద్ధాంతాన్ని
సూచిస్తుంది (తరువాత
చర్చించబోయే ఒకదాన్ని
\emph{తీసివేసిన} రూపం).
ఈ గౌరవానికి సెర్మెలో
అర్హుడు; ఎందుకంటే
\citeyear{Zermelo1908Untersuchungen}లో
ఈ సిద్ధాంతాన్ని ముఖ్యంగా
అతడే రూపొందించాడు.\footnote{చరిత్ర,
సాంకేతిక విషయాలపై
ఆసక్తికర వ్యాఖ్యల కోసం
\citet[Appendix
A]{Potter2004} చూడండి.}

ఈ సిద్ధాంతం అపారమైన
గణితాన్ని చేయగలిగేంత
శక్తివంతమైనది. ముఖ్యంగా,
మీరు \olref[sfr][][]{part}ను
తిరిగి \emph{చూసి}, అక్కడ
అనౌపచారికంగా చేసిన
ప్రతిదానినీ $\Zminus$లో
మరింత ఔపచారికంగా
చేయవచ్చని మీరే
నిర్ధారించుకోవాలి. (కొంతసేపు
అలా చేసిన తరువాత,
ముందుకు దాటి
\olref[sth][z][arbintersections]{sec}ను
చదవాలనుకోవచ్చు.)
కాబట్టి ఇకపై, మరో
వ్యాఖ్య లేకుండా, మనం
(కనీసం) $\Zminus$లో
పని చేస్తున్నామని
భావిస్తాం.

\end{document}
